\section{Архитектура система}

Решење описано у овом раду чине три компоненте које заједно симулирају и приказују једно паметно окружење: серверска апликација која игра улогу паметне зграде, веб контролна табла којом се током демонстрације мењају вредности сензора и \textit{Android} АР апликација као главни потрошач података. Ток података је једносмеран и једноставан: контролна табла шаље нове вредности сензора серверу, сервер их чува и обогаћује изведеним информацијама, а АР апликација их периодично преузима и исцртава у простору.

\begin{enumerate}
\item Презентер на контролној табли помери клизач или промени заузетост учионице; промена се после кратког задршка шаље серверу.
\item Сервер проверава да ли су вредности у дозвољеним границама, уписује их у меморијско стање и на сваки упит израчунава статусе и препоруку.
\item АР апликација, док прати маркер учионице, на сваке три секунде преузима стање те учионице и освежава панел усидрен изнад маркера.
\end{enumerate}

\subsection{Серверска апликација}

Сервер је \textit{Spring Boot} апликација (верзија 3.5, Java 21) која симулира паметно окружење. Стање учионица држи у меморији, у структури \texttt{ConcurrentHashMap}, без базе података: при покретању се учитавају три учионице (101, 102 и 103) са почетним вредностима температуре, буке и угљен-диоксида, заузетошћу и дневним распоредом часова. Целокупан изведени садржај рачуна се на серверу, па су оба клијента само прикази добијених података. РЕСТ интерфејс чине три крајње тачке приказане у табели~\ref{tab:endpoints}.

\begin{table}[H]
\centering
\small
\begin{tabular}{llp{5.6cm}}
\toprule
\textbf{Метода} & \textbf{Путања} & \textbf{Опис} \\
\midrule
\texttt{GET} & \texttt{/api/rooms} & листа свих учионица са тренутним стањем \\
\texttt{GET} & \texttt{/api/rooms/\{roomId\}} & стање једне учионице; 404 за непознат идентификатор \\
\texttt{PUT} & \texttt{/api/rooms/\{roomId\}/sensors} & упис нових вредности сензора и заузетости \\
\bottomrule
\end{tabular}
\caption{Крајње тачке РЕСТ интерфејса сервера}
\label{tab:endpoints}
\end{table}

За сваку измерену величину сервер израчунава статус са вредностима \texttt{OK}, \texttt{WARNING} или \texttt{CRITICAL} према задатим праговима: температура је у реду између 20 и 26~°C, бука испод 45~dB, а угљен-диоксид испод 800~ppm, док вредности изван ширих граница постају критичне. На основу статуса формира се и препорука на српском језику као подршка одлучивању --- од „Услови су оптимални." преко „Пратити услове у просторији." до „Хитно проветрити просторију." када је концентрација угљен-диоксида критична. Сервер из распореда часова и сопственог сата изводи и текст заузетости (назив текућег часа и време до којег је учионица заузета), а враћеним вредностима додаје благи синусни шум како би панели деловали живо и без сталних промена на контролној табли.

\subsection{Контролна табла}

Контролна табла је једнострана \textit{Angular} апликација (назив \texttt{kontrolna-tabla}) са по једном картицом за сваку учионицу. Картица садржи прекидач заузетости и три клизача --- температуру, буку и угљен-диоксид --- чија се свака промена, после задршка од 300~ms (енгл. \textit{debounce}), шаље серверу позивом \texttt{PUT /api/rooms/\{roomId\}/sensors}. Табла истовремено на сваких пет секунди преузима свеже стање свих учионица. Њена улога у систему је демонстрациона: пошто у учионицама не постоје стварни сензори, презентер преко клизача уживо симулира промене у паметном окружењу и посматра како се оне одражавају на АР панелу.

\snimakekrana{slike/01-kontrolna-tabla.png}{Контролна табла са картицама учионица и клизачима}

\subsection{АР апликација}

Трећа и централна компонента је \textit{Android} апликација заснована на \textit{ARCore} платформи, главни потрошач података из паметног окружења. Она препознаје маркер на вратима учионице, из његовог имена изводи идентификатор просторије, преузима њено стање са сервера и изнад маркера сидри панел са живим подацима. Развоју и имплементацији ове апликације посвећена су наредна два поглавља.
